\subsection{简单性定理}

引理 2. 设 H 是 G 的一个正规子群。存在 S 的一个子集 X，使得 BH = G\_X，并且 X 的每个元素都与 S-X 的每个元素交换。

由于 BH 是包含 B 的 G 的一个子群，定理 3 给出 S 的唯一子集 X，使得 \(BH = G_{X}\)。

设 \(s_1 \in X\) 且 \(s_2 \in S - X\)；设 \(n_1\) 和 \(n_2\) 分别是 \(N\) 中 \(s_1\) 和 \(s_2\) 的代表元。于是 \(n_1 \in G_X = BH\)，存在 \(b \in B\) 使得 \(bn_1 \in H\)。由于 \(H\) 是 \(G\) 的正规子群，元素 \(h = n_2bn_1n_2^{-1}\) 属于 \(G\)，且属于 \(H\)。这意味着

\[
h \in \mathrm{C} (s _ {2}). \mathrm{C} (s _ {1}). \mathrm{C} (s _ {2}).
\]

若 \(s_2s_1s_2\) 的长度等于 3，则定理 2 的推论 1 表明

\[
\mathrm{C} (s _ {2}). \mathrm{C} (s _ {1}). \mathrm{C} (s _ {2}) = \mathrm{C} (s _ {2} s _ {1} s _ {2}),
\]

从而 \(h \in H \cap C(s_{2}s_{1}s_{2})\)。由于 \(H \cap C(s_{2}s_{1}s_{2})\) 非空，\(s_{2}s_{1}s_{2} \in W_{X}\)。又因 \((s_{2}, s_{1}, s_{2})\) 是一个约化分解，故 \(s_{2} \in X\)，这与我们的假设矛盾。

因此 \(l_{\mathrm{S}}(s_2s_1s_2) \leqslant 2\)；若 \(l_{\mathrm{S}}(s_2s_1s_2) = 1\)，则 \(s_1s_2 \in \mathrm{S}\)，从而 \((s_1s_2)^2 = 1\)，或者 \(s_1s_2 = s_2s_1\)。若 \(l_{\mathrm{S}}(s_2s_1s_2) = 2\)，则第 1 节第 5 号的性质 (E) 蕴含 \(s_2s_1 = s_1s_2\)，因为 \(s_1 \neq s_2\)。证毕。

下面的群 U 的性质将在下文定理 5 中用到：

\begin{enumerate}
\def\labelenumi{(\Alph{enumi})}
\setcounter{enumi}{17}
\tightlist
\item
  对于 U 的任意不同于 U 的正规子群 V，U/V 的换位子群（参见《代数》第 I 章第 6 节第 8 号）不同于 U/V。
\end{enumerate}

每个可解群都满足 (R)；特别地，每个阿贝尔群都满足 (R)。可以证明，对称群 \(S_{n}\) 满足 (R)（参见习题 29）。

定理 5. 设 Z 为 B 的所有共轭子群的交，U 为 B 的一个子群，且 \(G_{1}\) 为 G 中由 U 的共轭子群生成的子群。作如下假设：

\begin{enumerate}
\def\labelenumi{\arabic{enumi})}
\item
  U 在 B 中正规且 B = UT。
\item
  U 具有性质 (R.)。
\item
  \(\mathbf{G}_1\) 等于其换位子群。
\item
  Coxeter 系统 (W, S) 是不可约的（参见~§ 1，第 9 号）。
\end{enumerate}

于是，G 中每个由 \(G_{1}\) 规范化的子群 H，或者包含于 Z，或者包含 \(G_{1}\)。

首先证明 \(G = G_{1}T\)。群 \(G_{1}T\) 包含 B，因而是自身的规范化子（定理 4）；但是由于 N 规范化 \(G_{1}\) 和 T，它也规范化 \(G_{1}T\)，故 \(N \subset G_{1}T\)。由于 G 由 B 和 N 生成，得到 \(G = G_{1}T\)。

接下来，置

\[
\begin{array}{c} \mathrm{G} ^ {\prime} = \mathrm{G} _ {1} \mathrm{H}, \quad \mathrm{B} ^ {\prime} = \mathrm{B} \cap \mathrm{G} ^ {\prime}, \quad \mathrm{N} ^ {\prime} = \mathrm{N} \cap \mathrm{G} ^ {\prime}, \\ \mathrm{T} ^ {\prime} = \mathrm{T} \cap \mathrm{G} ^ {\prime} = \mathrm{B} ^ {\prime} \cap \mathrm{N} ^ {\prime} \quad \text {and} \quad \mathrm{W} ^ {\prime} = \mathrm{N} ^ {\prime} / \mathrm{T} ^ {\prime}. \end{array}
\]

我们有 \(G = G'T\)，因为 \(G'\) 包含 \(G_{1}\)，因而 \(N = N'T\)。于是，\(N'\) 包含于 N 的嵌入通过取商定义了一个同构 \(\alpha : W' \to W\)。置 \(S' = \alpha^{-1}(S)\)。

现在证明 \((\mathrm{G}^{\prime},\mathrm{B}^{\prime},\mathrm{N}^{\prime},\mathrm{S}^{\prime})\) 是一个 Tits 系统。由于 G = BNB 且 B = TU = UT，有 G = UNU。由于 U 是 \(G^{\prime}\) 的一个子群，故 \(G^{\prime} = UN^{\prime}U\)；由于 \(U \subset B^{\prime}\)，这就证明了 (T1)。公理 (T2) 成立，因为 \(\alpha\) 是一个同构。设 \(w \in W\)，并令 \(w' = \alpha^{-1}(w)\) 为 \(W^{\prime}\) 中的相应元素。我们有

\[
\mathrm{B} w \mathrm{B} = \mathrm{B} w \mathrm{B} ^ {\prime} = \mathrm{B} w ^ {\prime} \mathrm{B} ^ {\prime}, \quad \text { since } \mathrm{B} = \mathrm{B} ^ {\prime} \mathrm{T}.
\]

由此得到 \(G' \cap BwB = B'w'B'\)，这意味着，\(G'\) 到 G 中的包含映射通过取商定义了一个从

\(B'\backslash G'/B'\) 到 \(B\backslash G/B\) 的双射。公理 (T3) 立即成立。公理 (T4) 由 \(B = B'T\) 得出。

子群 H 在 \(G'\) 中正规。将引理 2 应用于 \((G', B', N', S')\)，存在一个子集 \(X'\)，它是 \(S'\) 的子集，使得 \(B'H = G_{X'}'\)，且 \(S' - X'\) 的每个元素都与 \(X'\) 的每个元素交换。根据假设 (4)，只有两种可能：

\begin{enumerate}
\def\labelenumi{\alph{enumi})}
\item
  \(X' = \varnothing\)，即~\(B'H = B'\)，所以 \(H \subset B' \subset B\)。若 \(g \in G\)，则 \(g = g_1 t\)，其中 \(g_1 \in G_1\)、\(t \in T\)，且 \(H \subset g_1 Bg_1^{-1}\)，因为 \(G_1\) 规范化 \(H\)。因此 \(H \subset gBg^{-1}\)，并且由于 \(Z\) 是各个 \(gBg^{-1}\) 的交，我们有 \(H \subset Z\)。
\item
  \(X' = S'\)，即~\(B'H = G'\)。由于 \(G = G'T\)，我们有
\end{enumerate}

\[
\mathrm{G} = \mathrm{B} ^ {\prime} \mathrm{HT} = \mathrm{HB} ^ {\prime} \mathrm{T} = \mathrm{HB}.
\]

由于 B 规范化 U，U 的每个共轭子群都形如 \(hUh^{-1}\)，其中 \(h \in H\)。这样的子群包含于 UH 中，因而有 \(G_{1} \subset UH\)，根据 \(G_{1}\) 的定义。于是，我们有如下同构：

\[
\mathrm{U} / (\mathrm{U} \cap \mathrm{H}) \cong \mathrm{UH} / \mathrm{H} = \mathrm{G} _ {1} \mathrm{H} / \mathrm{H} \cong \mathrm{G} _ {1} / (\mathrm{G} _ {1} \cap \mathrm{H}).
\]

根据假设 (3)，\(G_{1}/(G_{1} \cap H)\) 等于其换位子群。假设 (2) 现在表明，群 \(U/(U \cap H)\)（它同构于 \(G_{1}/(G_{1} \cap H)\)）退化为单位元。因此 \(G_{1} \cap H = G_{1}\) 且 \(G_{1} \subset H\)，这就完成了证明。

推论。在定理 5 的假设下，群 \(\mathrm{G}_{1}/(\mathrm{G}_{1} \cap \mathrm{Z})\) 或者是非阿贝尔单群，或者退化为单位元。

定理 5 表明，\(G_{1}/(G_{1} \cap Z)\) 是单群或退化为单位元。另一方面，假设 (3) 蕴含它等于其换位子群。因此推论得证。

备注。1) 证明 \((\mathbf{G}', \mathbf{B}', \mathbf{N}', \mathbf{S}')\) 是一个 Tits 系统时，并未使用假设 (2)、(3)、(4)。

\begin{enumerate}
\def\labelenumi{\arabic{enumi})}
\setcounter{enumi}{1}
\item
  假设 \(\mathbf{Z} \cap \mathbf{U} = \{1\}\)。由于 \(\mathbf{Z}\) 和 \(\mathbf{U}\) 在 \(\mathbf{B}\) 中正规，故 \(\mathbf{Z}\) 的每个元素都与 \(\mathbf{U}\) 的每个元素交换，因而也与 \(\mathbf{G}_1\) 的每个元素交换。根据前一推论，\(\mathbf{G}_1 \cap \mathbf{Z}\) 是 \(\mathbf{G}_1\) 的中心。
\item
  假设 (3) 由下述条件蕴含：
\end{enumerate}

(3') U 由换位子 \(b^{-1}u^{-1}bu\) 生成，其中 \(u \in \mathbf{U}\) 且 \(b \in \mathbf{B} \cap \mathbf{G}_1\)。

例子。1) 设 \(k\) 是一个域，\(n\) 是满足 \(\geqslant 0\) 的整数，\(G = \mathbf{GL}(n, k)\)，并设 \((G, B, N, S)\) 是第 2 号中描述的 Tits 系统。令 \(U\) 为 \(G\) 的严格上三角子群，即 \(B\) 中由对角线元素等于 1 的矩阵组成的子群。定理 5 中的条件 (1) 是直接的，条件 (2) 也直接成立，因为 U 是可解的。当 \(n \geqslant 2\) 时，条件 (4) 得到满足。可以证明（参见~《代数》第 II 章第 10 节习题 13），当 \(n \geqslant 3\) 且 \(\operatorname{Card}(k) \geqslant 4\) 时，条件 (3) 得到满足。在这些条件下，我们得到 \(G_{1}/(G_{1} \cap Z)\) 是单群，并且 \(G_{1} \cap Z\) 是 \(G_{1}\) 的中心（参见~备注 2）。

当 \(k\) 为交换域时，\(G_1 = \mathbf{SL}(n,k)\)（参见~《代数》第 III 章第 8 节第 9 号）。

*2) 设 \(\mathfrak{g}\) 是 \(\mathbf{C}\) 上的一个单李代数，并设 \(\mathbf{G}\) 是 \(\mathfrak{g}\) 的内自同构群（参见~第 III 章第 6 节第 2 号，命题 2）。利用定理 5 可以证明，\(\mathbf{G}\) 是非阿贝尔单群。*

